\documentclass[../../../include/open-logic-section]{subfiles}

\begin{document}

\olfileid{sth}{card-arithmetic}{ch}

\olsection{सांतत्य गृहीतक}

मागील निष्कर्ष योग्यच सूचित करतो की अनंत संचांकांचा $2$ च्या घातांशी, म्हणजे (\olref[opps]{lem:SizePowerset2Exp} नुसार) घातसंच घेण्याशी, संबंध सरळ असेल, तर संचांक घातांकन अगदी \emph{सोपे} होईल. पण अनंत संचांकांचा घातसंचांशी असा सरळ संबंध \emph{असतोच}, असे आपण गृहीत धरू शकत नाही.

हे उलगडण्यास सुरुवात करण्यासाठी काही सोयीची संकेतरचना मांडू.

\begin{defn}
$\cardsucc{\cardfont{a}}$ हा $\cardfont{a}$ पेक्षा काटेकोरपणे मोठा असलेला लघुतम संचांक, म्हणजे त्याचा उत्तरवर्ती संचांक, असू द्या. पुढील दोन अनंत अनुक्रम परिभाषित करू:
\begin{align*}
	\aleph_{0} &\defis \omega &
	\beth_{0} &\defis \omega\\
	\aleph_{\alpha \ordplus 1} &\defis \cardsucc{(\aleph_{\alpha})} &
	\beth_{\alpha+1} &\defis \cardexpo{2}{\beth_{\alpha}}\\
	\aleph_{\alpha} &\defis \bigcup_{\beta< \alpha} \aleph_{\beta} &
	\beth_{\alpha} &\defis \bigcup_{\beta < \alpha}\beth_{\beta} & \text{जेव्हा $\alpha$ सीमा क्रमसूचक संख्या असते}.
	\end{align*}
\end{defn}

$\cardsucc{\cardfont{a}}$ ची व्याख्या योग्य आहे. कारण \olref[cardinals][classing]{lem:NoLargestCardinal} नुसार प्रत्येक संचांक $\cardfont{a}$ साठी $\cardfont{a}$ पेक्षा मोठा संचांक असतो, आणि अनंतपार विगमनामुळे $\cardfont{a}$ पेक्षा मोठा असा \emph{लघुतम} संचांक असल्याची खात्री मिळते. $\cardfont{a}$ यांसारखे संचांक देणाऱ्या या दोन अनुक्रमांची उर्वरित व्याख्या अनंतपार पुनरावर्तनाने मिळते.

कँटरने “$\aleph$” ही संकेतरचना आणली. हे \emph{अलेफ} आहे: हिब्रू वर्णमालेतील पहिले अक्षर आणि हिब्रू भाषेतील “अनंत” या शब्दाचेही पहिले अक्षर. पर्सने “$\beth$” ही संकेतरचना आणली. हे \emph{बेथ} आहे: हिब्रू वर्णमालेतील दुसरे अक्षर.\footnote{डिसेंबर 1900 मध्ये कँटरला लिहिलेल्या पत्रात पर्सने ही संकेतरचना वापरली. दुर्दैवाने त्याच पत्रात $\alpha \geq \omega$ साठी $\beth_\alpha$ अस्तित्वात नसल्याचा चुकीचा युक्तिवादही त्याने दिला.} या संकेतरचना आपल्याला अनंत संचांक देतात.

\begin{prop}
प्रत्येक क्रमसूचक संख्या $\alpha$ साठी $\aleph_\alpha$ आणि $\beth_\alpha$ हे संचांक आहेत.
\end{prop}

\begin{proof}
दोन्ही निष्कर्ष साध्या अनंतपार विगमनाने मिळतात. \olref[cardinals][classing]{omegaisacardinal} नुसार $\aleph_0 =
\beth_0 = \omega$ हा संचांक आहे. $\aleph_\alpha$ आणि $\beth_\alpha$ हे दोन्ही संचांक आहेत असे गृहीत धरल्यास, $\aleph_{\alpha+1}$ आणि $\beth_{\alpha+1}$ हे व्याख्येनुसारच संचांक आहेत. शिवाय \olref[cardinals][classing]{unioncardinalscardinal} नुसार संचांकांच्या संचाचे युतीकरण हा संचांक असतो.
\end{proof}
\noindent
तसेच प्रत्येक अनंत संचांक हा $\aleph$ च्या स्वरूपाचा असतो:

\begin{prop}
$\cardfont{a}$ हा अनंत संचांक असेल, तर एकमेव $\gamma$ साठी $\cardfont{a} =
\aleph_\gamma$.
\end{prop}

\begin{proof}
अनंत संचांकांवर अनंतपार विगमन वापरू. अलेफ-शून्याच्या व्याख्येमुळे लघुतम अनंत संचांकाचे आधारप्रकरण मिळते. पुढील विगमनात $\cardfont{b} < \cardfont{a}$ असलेल्या प्रत्येक अनंत संचांकासाठी $\cardfont{b} =
\aleph_{\gamma_\cardfont{b}}$ असे गृहीत धरा; खालील युतीकरणांतही अशाच अनंत संचांकांचा विचार केला आहे. एखाद्या $\cardfont{b}$ साठी $\cardfont{a} =
\cardsucc{\cardfont{b}}$ असल्यास $\cardfont{a} =
\cardsucc{(\aleph_{\gamma_\cardfont{b}})}=
\aleph_{\gamma_\cardfont{b}+1}$. जर $\cardfont{a}$ हा कोणत्याही संचांकाचा उत्तरवर्ती नसेल, तर संचांक या क्रमसूचक संख्या असल्यामुळे $\cardfont{a} = \bigcup_{\cardfont{b} < \cardfont{a}} \cardfont{b} =
\bigcup_{\cardfont{b} < \cardfont{a}}{\aleph_{\gamma_\cardfont{b}}}$.
म्हणून $\gamma =
\bigcup_{\cardfont{b} < \cardfont{a}}\gamma_\cardfont{b}$ घेतल्यास $\cardfont{a} = \aleph_\gamma$. अलेफ अनुक्रम काटेकोरपणे वाढतो: उत्तरवर्ती टप्प्यावर मोठा संचांक घेतला जातो, आणि सीमाटप्प्यावर आधीच्या सर्व मूल्यांचे युतीकरण केले जाते. त्यामुळे भिन्न निर्देशांकांना एकच मूल्य मिळू शकत नाही; म्हणून निर्देशांक एकमेव आहे.
\end{proof}

प्रत्येक अनंत संचांक हा $\aleph$ च्या स्वरूपाचा असल्यामुळे पुढील प्रश्न पडतो: प्रत्येक अनंत संचांक हा $\beth$ च्याही स्वरूपाचा असतो का? असे \emph{असते}, तर अनंत संचांकांचा घातसंच घेण्याच्या क्रियेशी सरळ संबंध असता. खरे तर पुढील विधान मिळाले असते:

\begin{defish}
\emph{सामान्यीकृत सांतत्य गृहीतक} (GCH). प्रत्येक $\alpha$ साठी $\aleph_\alpha = \beth_\alpha$.
\end{defish}

तसेच GCH खरे असेल, तर संचांक घातांकनात बरीच सुलभीकरणे करता येतील. विशेषतः अनंत आधार आणि शून्येतर घातांकांच्या बाबतीत, $\cardfont{b} < \cardfont{a}$ असताना $\cardexpo{\cardfont{a}}{\cardfont{b}}$ चे मूल्य
$\cardfont{a}\leq \cardexpo{\cardfont{a}}{\cardfont{b}} \leq
\cardsucc{\cardfont{a}}$ या मर्यादांत असल्याचे दाखवता येईल. त्यानंतर या दोन शक्यतांपैकी कोणती घडते (म्हणजे $\cardfont{a} = \cardexpo{\cardfont{a}}{\cardfont{b}}$ की $\cardexpo{\cardfont{a}}{\cardfont{b}} =
\cardsucc{\cardfont{a}}$), हे ठरवणाऱ्या नेमक्या अटीही देता येतील.\footnote{ही अट \emph{सहअंतिमतेवर} ठरते.}

पण GCH हे \emph{गृहीतक} आहे, \emph{प्रमेय} नाही. खरे तर \citet{Godel1938} यांनी असे सिद्ध केले की $\ZFC$ सुसंगत असेल, तर $\ZFC + \text{GCH}$ सुद्धा सुसंगत आहे. पण नंतर असेही आढळले की $\ZFC$ मध्ये $\lnot$GCH तितक्याच रीतीने जोडता येते. यासाठी GCH चे सर्वांत सोपे, क्षुल्लक नसलेले \emph{विशेषप्रकरण} पाहा:

\begin{defish}
\emph{सांतत्य गृहीतक} (CH). $\aleph_1 = \beth_1$.
\end{defish}

\citet{Cohen1963} यांनी असे सिद्ध केले की $\ZFC$ सुसंगत असेल, तर $\ZFC
+ \lnot\text{CH}$ सुद्धा सुसंगत आहे. म्हणून सांतत्य गृहीतक हे $\ZFC$ पासून स्वतंत्र आहे.

सांतत्य गृहीतकाला हे नाव पडण्याचे कारण असे की “सांतत्य” हे वास्तव संख्यारेषा $\Real$ हिचे दुसरे नाव आहे. \olref[opps]{continuumis2aleph0} नुसार $\card{\Real} =
\beth_1$. म्हणून सांतत्य गृहीतक असे सांगते की नैसर्गिक संख्यांचा संचांक $\aleph_0 = \beth_0$ आणि सांतत्याचा संचांक $\beth_1$ यांच्या दरम्यान कोणताही संचांक नाही.

(G)CH हे $\ZFC$ पासून \emph{स्वतंत्र} असल्यामुळे त्यांच्या \emph{सत्यतेविषयी} काय म्हणावे? याविषयी \emph{बरेच} काही सांगता येते. संच उपपत्तीतील अलीकडील अनेक फलदायी संशोधनांचा रोख या प्रश्नांच्या तपासाकडे आहे. पण येथे दोन मुद्दे थोडक्यात अधोरेखित करू.

पहिला मुद्दा: या औपचारिक स्वातंत्र्यनिष्कर्षांवरून GCH किंवा CH यांचे सत्यमूल्य \emph{अनिर्धारित} आहे, असा निष्कर्ष \emph{लगेच} मिळत नाही. अखेर, आपल्याला स्वाभाविक वाटणारी आणखी स्वयंसिद्धके जोडल्यावर हा प्रश्न एका किंवा दुसऱ्या बाजूने सुटू शकेल. खुद्द गोडेल यांनीही हा योग्य प्रतिसाद असल्याचे सुचवले होते.

दुसरा मुद्दा: CH हे $\ZFC$ पासून स्वतंत्र असणे नक्कीच \emph{लक्षणीय} आहे; पण ते शब्दशः \emph{अविश्वसनीय} नाही. मुद्दा एवढाच की $\ZFC$ आपल्याला जे सांगते त्यानुसार संचांकापासून त्याच्या उत्तरवर्ती संचांकाकडे जाण्यासाठी, केवळ घातसंच घेण्यापेक्षा अधिक सूक्ष्म साधन लागू शकते.
 %घातसंच घेण्याची क्रिया श्रेणीच्या एका टप्प्यापासून त्याच्या
 %उत्तरवर्ती टप्प्याकडे नेते (म्हणजे $V_{\alpha}$ पासून
 %$V_{\alpha+1}$ कडे).

हे दोन मुद्दे मांडल्यानंतर, आणखी जाणून घ्यायचे असल्यास या चर्चेत आपापली भूमिका मांडणाऱ्या विविध तत्त्वज्ञ आणि गणितज्ञांच्या लेखनाचा आधार घ्यावा लागेल.\footnote{तरी सुरुवात \citet[\S15.6]{Potter2004} यांच्या विवेचनापासून करणे उपयुक्त ठरेल.}

\end{document}
